\chapter{Moving packets, dynamic frames, and sparse transfer graphs}
\label{ch:dynamic-frames}
\index{moving packet}
\index{time-dependent projector}
\index{packet transfer graph}

\section{Why one fixed projector is not the end of the theory}
The exact packet theorem in \cref{ch:exact-simulation} follows a state inside one fixed subspace. That is sufficient for a conserved charge, but it is not necessary for classical collapse. A circuit may move through a sequence of small subspaces, or it may preserve a small packet only after a time-dependent change of coordinates. The correct object is therefore a \emph{packet schedule}, not necessarily one static projector.

\begin{definition}[Packet schedule]
For a circuit
\[
U=G_m\cdots G_1
\]
on a Hilbert space $\Hh$, a packet schedule is a sequence of orthogonal projectors
\[
P_0,P_1,\ldots,P_m.
\]
It is exact when
\begin{equation}
G_tP_{t-1}=P_tG_tP_{t-1}
\qquad(1\le t\le m).
\label{eq:dynamic-packet-condition}
\end{equation}
\end{definition}

The condition says that gate $G_t$ maps the currently active packet into the next active packet. The dimensions may vary with time.

\section{Exact dynamic-packet simulation}
Choose isometries
\[
B_t:\C^{D_t}\longrightarrow\Hh,
\qquad
B_tB_t^*=P_t,
\qquad
B_t^*B_t=I_{D_t},
\]
where $D_t=\rank P_t$. Define the rectangular coordinate update
\begin{equation}
M_t=B_t^*G_tB_{t-1}.
\label{eq:dynamic-coordinate-update}
\end{equation}

\begin{theorem}[Exact dynamic-packet theorem]
\label{thm:dynamic-packet}
Suppose $P_0\ket{\psi_0}=\ket{\psi_0}$ and \cref{eq:dynamic-packet-condition} holds for every gate. Put
\[
z_0=B_0^*\ket{\psi_0},
\qquad
z_t=M_tz_{t-1}.
\]
Then
\[
G_t\cdots G_1\ket{\psi_0}=B_tz_t
\]
for every $t$.
\end{theorem}

\begin{proof}
The claim is true at $t=0$. If the state after $t-1$ gates is $B_{t-1}z_{t-1}$, then
\[
G_tB_{t-1}z_{t-1}
=P_tG_tP_{t-1}B_{t-1}z_{t-1}
=B_tB_t^*G_tB_{t-1}z_{t-1}
=B_tM_tz_{t-1}.
\]
This is $B_tz_t$.
\end{proof}

\begin{theorem}[Polynomial moving-packet collapse]
\label{thm:moving-packet-collapse}
Let $\{C_x\}$ be a polynomial-time uniform circuit family. Suppose an exact packet schedule can be produced from $x$ such that:
\begin{enumerate}
\item $\max_tD_t\le\poly(|x|)$;
\item the coordinate description of the initial state is computable in polynomial time;
\item each rectangular update $M_t$ can be applied to polynomial precision in polynomial time;
\item the final acceptance probability is computable from $z_m$ in polynomial time.
\end{enumerate}
Then the language decided by the circuit family belongs to $P$.
\end{theorem}

\begin{proof}
The simulator stores at most polynomially many coordinates at every step and applies polynomially many polynomial-time updates. \Cref{thm:dynamic-packet} proves exactness before numerical rounding. Polynomial precision below the amplified promise gap gives a deterministic polynomial-time decision procedure.
\end{proof}

\begin{derivedbox}
A circuit does not need one globally conserved sector. It is enough that the state follow an efficiently discoverable path through polynomial-dimensional sectors.
\end{derivedbox}

\section{Moving frames and hidden symmetry}
The packet schedule has a useful coordinate form. Let $P$ be one fixed projector and let $W_t$ be unitaries with
\[
P_t=W_tPW_t^*.
\]
In the moving frame, the transformed gate is
\begin{equation}
\widetilde G_t
=W_t^*G_tW_{t-1}.
\label{eq:moving-frame-gate}
\end{equation}
The schedule condition is equivalent to
\[
\widetilde G_tP=P\widetilde G_tP.
\]

\begin{corollary}[Moving-frame collapse]
\label{cor:moving-frame-collapse}
Suppose a polynomial-time uniform circuit admits efficiently described frames $W_t$ for which the transformed gates in \cref{eq:moving-frame-gate} preserve one fixed polynomial-dimensional packet, and their restricted matrices and the transformed final measurement are efficiently computable. Then the circuit family is classically simulable in polynomial time.
\end{corollary}

This is the appropriate theorem for a symmetry that changes orientation with time. Examples to test include Clifford frames, Gaussian mode frames, rotating spin axes, and gauge-fixed local coordinates.

\begin{scopebox}
The frame must be independently discoverable and efficiently usable. Choosing $W_t=G_t\cdots G_1$ makes the transformed state trivial but is circular: computing that frame is the original simulation problem. The theorem is useful only when the frames have a compact classical description and their restricted actions can be computed without constructing the full state.
\end{scopebox}

\begin{figure}[H]
\centering
\includegraphics[width=.94\textwidth]{moving_frame_packets.pdf}
\caption{A circuit may fail to preserve one laboratory-frame packet while remaining inside a small packet in a computable moving frame.}
\end{figure}

\section{Sparse packet-transfer graphs}
Let
\[
I=\sum_{\alpha\in\Lambda}P_\alpha
\]
be an orthogonal packet decomposition. Write the gate blocks as
\[
G_{t;\beta\alpha}=P_\beta G_tP_\alpha.
\]
Starting from an active packet set $S_0$, define recursively
\begin{equation}
S_t
=
\left\{
\beta:\exists\alpha\in S_{t-1}
\text{ with }G_{t;\beta\alpha}\ne0
\right\}.
\label{eq:reachable-packet-set}
\end{equation}

\begin{corollary}[Reachable packet-graph collapse]
\label{cor:packet-graph-collapse}
Suppose the packet sets $S_t$ and nonzero transfer blocks are computable in polynomial time, and
\[
D_t
=
\sum_{\alpha\in S_t}\rank P_\alpha
\le\poly(n)
\]
for every $t$. If all active block updates and the final measurement are efficiently computable, then the circuit family is classically polynomial-time simulable.
\end{corollary}

\begin{proof}
Take
\[
P_t=\sum_{\alpha\in S_t}P_\alpha.
\]
The definition of $S_t$ gives $G_tP_{t-1}=P_tG_tP_{t-1}$. Apply \cref{thm:moving-packet-collapse}.
\end{proof}

This theorem permits exact sector changes. What matters is not that one charge be fixed, but that the union of reachable sectors remain polynomial.

\begin{corollary}[Bounded-excitation collapse]
\label{cor:bounded-excitation-collapse}
Consider $n$ qubits and suppose the initial state and every circuit prefix are certified to lie in
\[
\Hh_{\le K}
=
\bigoplus_{k=0}^{K}\Hh_{n,k}
\]
for a fixed constant $K$. If the gates and measurement have polynomial-time actions in the occupation basis, then the circuit family lies in $P$, even when individual gates change excitation number.
\end{corollary}

\begin{proof}
The active dimension is
\[
\sum_{k=0}^{K}\binom nk
=O(n^K).
\]
Use the weight sectors as packet nodes and apply \cref{cor:packet-graph-collapse}.
\end{proof}

For example, a fixed-weight input followed by only a constant number of pair-creation or pair-annihilation events remains in a constant collection of low-weight sectors. Exact number conservation is therefore stronger than necessary.

\begin{figure}[H]
\centering
\includegraphics[width=.90\textwidth]{reachable_packet_graph.pdf}
\caption{Sector-changing gates remain classically tractable when the reachable packet graph and total active dimension stay polynomial.}
\end{figure}

\section{Approximate dynamic packets}
Define the one-step leakage from the scheduled packet by
\[
\delta_t
=
\|(I-P_t)G_tP_{t-1}\|.
\]
Let
\[
\widehat G_t=P_tG_tP_{t-1}.
\]

\begin{theorem}[Dynamic sum-of-leakages bound]
\label{thm:dynamic-leakage}
For an initial unit vector in $P_0\Hh$,
\[
\left\|
G_m\cdots G_1\ket{\psi_0}
-
\widehat G_m\cdots\widehat G_1\ket{\psi_0}
\right\|
\le
\sum_{t=1}^{m}\delta_t.
\]
For any acceptance operator $0\le M\le I$, the corresponding acceptance probabilities differ by at most
\[
2\sum_{t=1}^{m}\delta_t.
\]
\end{theorem}

\begin{proof}
A telescoping expansion inserts $G_t-\widehat G_t$ at one time. On $P_{t-1}\Hh$ this difference equals $(I-P_t)G_tP_{t-1}$, while all other unitary or projected factors have norm at most one. Summing the operator norms proves the state bound. If $\psi$ and $\phi$ have norm at most one, then
\[
|\langle\psi,M\psi\rangle-\langle\phi,M\phi\rangle|
\le
(\|\psi\|+\|\phi\|)\|\psi-\phi\|
\le2\|\psi-\phi\|.
\]
\end{proof}

\begin{corollary}[Approximate moving-packet collapse]
If the scheduled dimensions and updates are polynomial and
\[
\sum_t\delta_t\le\frac1{\poly(n)},
\]
then, after ordinary promise-gap amplification, the circuit family remains deterministically classically simulable in polynomial time.
\end{corollary}

The result is conservative but fully composable. It extends the fixed-projector estimate of \cref{ch:approximate-symmetry} to moving frames and changing packet sets.
